Coding agents open many pull requests that claim to make software faster.
Prior studies read these claims or re-execute a few of them; a stopwatch shows
whether code got faster, but not whether it got faster for the reason the
agent gives. We check agents' performance claims at the kernel level. From
\nRuntimePRs runtime-performance pull requests in the AIDev dataset we
extract the mechanism each one claims, map it to the kernel and hardware
signals that must move if the claim is true, rebuild both versions with a
build agent that never sees the claim, and measure them with untraced
timing, hardware counters and LTTng kernel traces. Agents rarely back their
claims: \pctReasonNone offer only reasoning or no evidence, and
\pctEffect state any effect size. In first measurements of \nAnalysed PRs,
\nFaster ran faster, \nNoChange showed no clear change and \nSlower ran
slower; the claimed mechanism left its kernel footprint in \nHeld, and
\nBehaviour also changed their output. Kernel signals separate claimed
from unclaimed mechanisms far better than the instruction count
(\kRatioFast$\times$ against \iRatioFast$\times$ for faster PRs), showing
what looking below the stopwatch adds.
